\usepackage[a4paper, margin=2.5cm]{geometry}

\usepackage{amsmath,amssymb,amsthm,mathtools}

% Em-dash glyph swapped to en-dash for narrower rendered dashes
\usepackage{fontspec}
\directlua{
  fonts.handlers.otf.addfeature {
    name = "shrinkemdash",
    type = "substitution",
    data = { [0x2014] = 0x2013 },
  }
}
\setmainfont{Latin Modern Roman}[RawFeature={+shrinkemdash}]
\setsansfont{Latin Modern Sans}[RawFeature={+shrinkemdash}]
\setmonofont{Latin Modern Mono}[RawFeature={+shrinkemdash}]

% Save \sqsubseteq before semantic redefines it (needed for tikz-cd labels)
\let\sqsubseteqOriginal\sqsubseteq
\let\sqsupseteqOriginal\sqsupseteq

\usepackage{semantic}        % \inference
\usepackage{stmaryrd}        % \llbracket, \rrparenthesis, ...
\newcommand{\Circle}{\bigcirc} % wasysym incompatible with lualatex

\usepackage[dvipsnames]{xcolor}
\usepackage{soul}             % \hl

\usepackage{booktabs}
\usepackage{tabularx}
\usepackage{arydshln}         % dashed col separators

\usepackage{graphicx}
\usepackage{subcaption}
\usepackage[labelfont=bf]{caption}

\usepackage{tikz-cd}
\usetikzlibrary{positioning,calc,fit,backgrounds,decorations.pathmorphing,shapes.geometric}
\tikzcdset{
  arrows={->},
  diagrams={column sep=large, row sep=large}
}

\usepackage[most]{tcolorbox}
\usepackage{fancyhdr}
\usepackage[normalem]{ulem}

\usepackage[backend=biber,style=numeric,sorting=nyt]{biblatex}
\addbibresource{refs.bib}

\usepackage{hyperref}
\hypersetup{
  colorlinks=true,
  linkcolor=blue,
  citecolor=magenta,
  urlcolor=cyan,
  pdftitle={Polymorphic Type Slicing},
  pdfauthor={Max Carroll},
}
\usepackage{cleveref}

\title{Polymorphic Type Slicing}
\author{Max Carroll \\ Sidney Sussex College \\ University of Cambridge}
\date{}

\definecolor{commentgreen}{RGB}{2,112,10}
\definecolor{weborange}{RGB}{255,165,0}
\definecolor{frenchplum}{RGB}{129,20,83}

% Vibrant overrides for dvipsnames slice/highlight colours
\definecolor{BrickRed}{RGB}{210,35,40}
\definecolor{NavyBlue}{RGB}{30,80,225}
\definecolor{OliveGreen}{RGB}{35,165,45}
\definecolor{Mulberry}{RGB}{215,55,175}

% Universal-property diagram convention: assumed (Class 2), induced (Class 3)
\colorlet{assumed}{BurntOrange!80!black}
\colorlet{induced}{TealBlue!45!black}

% Lattice ops and precision in a muted colour
\definecolor{latticecolor}{HTML}{8055B5}
\let\OLDsqcap\sqcap
\let\OLDsqcup\sqcup
\let\OLDsqsubset\sqsubset
\let\OLDsqsupset\sqsupset
\def\sqcap{\mathbin{{\color{latticecolor}\OLDsqcap}}}
\def\sqcup{\mathbin{{\color{latticecolor}\OLDsqcup}}}
\def\sqsubseteq{\mathrel{{\color{latticecolor}\sqsubseteqOriginal}}}
\def\sqsupseteq{\mathrel{{\color{latticecolor}\sqsupseteqOriginal}}}
\def\sqsubset{\mathrel{{\color{latticecolor}\OLDsqsubset}}}
\def\sqsupset{\mathrel{{\color{latticecolor}\OLDsqsupset}}}
\newcommand{\heytingto}{\mathbin{{\color{latticecolor}\to}}}
\newcommand{\heytingminus}{\mathbin{{\color{latticecolor}\setminus}}}
\newcommand{\latticetop}{{\color{latticecolor}\top}}
\newcommand{\latticebot}{{\color{latticecolor}\bot}}

% Highlighting
\newcommand{\hlc}[2][yellow]{{%
    \colorlet{foo}{#1}%
    \sethlcolor{foo}\hl{#2}}%
}
\newcommand{\error}[1]{\hlc[pink!50]{#1}}
\newcommand{\hlcmaths}[2][yellow!50]{\colorbox{#1}{$\displaystyle #2$}}

\newcommand{\dyn}{{\color{teal}\texttt{?}}}          % dynamic type
\newcommand{\gap}{{\color{teal}\Box}}                 % gap □
\newcommand{\hole}[1][]{{\color{teal}\llparenthesis} #1 {\color{teal}\rrparenthesis}} % hole ⦇·⦈
\newcommand{\cmark}{{\color{blue}\Circle}}            % context mark
\newcommand{\type}[1]{\llbracket #1 \rrbracket}       % type interpretation
\newcommand{\id}{\mathrm{id}}
\newcommand{\code}[1]{\texttt{#1}}

% Well-formedness: Δ ⊢ τ wf
\newcommand{\wf}[2][\Delta]{#1 \vdash #2\;\mathsf{wf}}

% Synthesis: Δ; Γ ⊢ e ⇑ τ;   #1 = ctx (default Δ;\,Γ), #2 = e, #3 = τ
\newcommand{\synthesis}[3][\Delta;\,\Gamma]{#1 \vdash #2 {\color{BrickRed}\ \Uparrow}\ #3}
% Analysis: Δ; Γ ⊢ e ⇓ τ
\newcommand{\analysis}[3][\Delta;\,\Gamma]{#1 \vdash #2 {\color{BlueViolet}\ \Downarrow}\ #3}

% Rule-name labels (text mode, for \texttt compatibility)
\newcommand{\synr}[1]{\textnormal{{\color{BrickRed}$\Uparrow$}\texttt{#1}}}
\newcommand{\anar}[1]{\textnormal{{\color{BlueViolet}$\Downarrow$}\texttt{#1}}}
\newcommand{\cnsr}[1]{\textnormal{$#1$}}

% Type assignment (internal language): Δ;Γ ⊢ d : τ
\newcommand{\typeassignment}[3][\Delta;\Gamma]{#1 \vdash #2 : #3}

% Synthesis slice: Γ ⊢ e ⟹ τ | S
\newcommand{\synthesisslice}[4][\Gamma]{#1 \vdash #2 {\color{BrickRed}\ \Rightarrow}\ #3 \mid #4}
% Analysis slice: Γ; C ⊢ e ⟸ τ | S
\newcommand{\analysisslice}[5][\Gamma]{#1; #2 \vdash #3 {\color{BlueViolet}\ \Leftarrow}\ #4 \mid #5}

% Casts: safe, error, chained
\newcommand{\scast}[2]{{\color{blue}\langle} #1 {\color{blue}\Rightarrow} #2 {\color{blue}\rangle}}
\newcommand{\scasterror}[2]{{\color{red}\langle} #1 {\color{red}\Rightarrow} \dyn {\color{red}\nRightarrow} #2 {\color{red}\rangle}}
\newcommand{\scastcast}[3]{{\color{blue}\langle} #1 {\color{blue}\Rightarrow} #2 {\color{blue}\Rightarrow} #3 {\color{blue}\rangle}}

% Type matching: ▸→, ▸ground, ▸∀
\newcommand{\funmatch}{\blacktriangleright_\rightarrow}
\newcommand{\groundmatch}{\blacktriangleright_{\text{ground}}}
\newcommand{\forallmatch}{\blacktriangleright_\forall}

% Elaboration judgements (with/without slice columns)
\newcommand{\elaborationSynthesis}[5][\Gamma]{#1 \vdash #2 {\color{BrickRed}\ \Uparrow\ } #3 \leadsto #4 \dashv #5}
\newcommand{\elaborationAnalysis}[6][\Gamma]{#1 \vdash #2 {\color{BlueViolet}\ \Downarrow\ } #3 \leadsto #4 : #5 \dashv #6}
\newcommand{\elaborationSynthesisSlice}[6][\Gamma]{#1 \vdash #2 {\color{BrickRed}\ \Rightarrow\ } #3 \mid #4 \leadsto #5 \dashv #6}
\newcommand{\elaborationAnalysisSlice}[9][\Gamma]{#1;\ #2 \vdash #3 {\color{BlueViolet}\ \Leftarrow\ } #4 \mid #5 \leadsto #6 : #7 \mid #8 \dashv #9}

% Context classification: Δ; Γ ⊢ C at d ▷ Δ'; Γ' [ m ]
\newcommand{\ctxclass}[5][\Delta;\,\Gamma]{#1 \vdash #2\;\mathbf{at}\;#3\;\triangleright\;#4\;[\,#5\,]}

% Marking arrow ↬; \marks is a TeX primitive so we use \marksto
\newcommand{\marksto}{\looparrowright}

% Synthesis/analysis marking judgements: Δ; Γ ⊢ e ↬ ě ⇑/⇓ τ
\newcommand{\synmark}[4][\Delta;\,\Gamma]{#1 \vdash #2 \marksto #3 {\color{BrickRed}\ \Uparrow}\ #4}
\newcommand{\anamark}[4][\Delta;\,\Gamma]{#1 \vdash #2 \marksto #3 {\color{BlueViolet}\ \Downarrow}\ #4}

% Marked context classification: Δ; Γ ⊢ C ↬ Č at d ▷ Δ'; Γ' [ m ]
\newcommand{\mctxclass}[6][\Delta;\,\Gamma]{#1 \vdash #2 \marksto #3\;\mathbf{at}\;#4\;\triangleright\;#5\;[\,#6\,]}

% Marked-expression accent ě (\check shadowed by LaTeX's \check)
\newcommand{\chk}[1]{\check{#1}}

% Mark wrappers ⦅·⦆ with a subscript glyph
\definecolor{markcolor}{RGB}{170,40,60}
\newcommand{\markopen}{{\color{markcolor}\llparenthesis}}
\newcommand{\markclose}{{\color{markcolor}\rrparenthesis}}

% Type-inconsistency mark ⦅ě⦆_{τ_syn ≁ τ_ana}
\newcommand{\markincon}[3][\cdot]{{\markopen #1 \markclose}_{#2 \not\sim #3}}

% Head-shape marks (focus's syn type fails join with □⇒□ / □+□ / □×□ / ∀·□)
\newcommand{\markarrow}[1][\cdot]{{\markopen #1 \markclose}_{\blacktriangleright \Rightarrow}}
\newcommand{\marksum}[1][\cdot]{{\markopen #1 \markclose}_{\blacktriangleright +}}
\newcommand{\markprod}[1][\cdot]{{\markopen #1 \markclose}_{\blacktriangleright \times}}
\newcommand{\markforall}[1][\cdot]{{\markopen #1 \markclose}_{\blacktriangleright \forall}}

% Syntactic-limitation marks (bare λ⇒ / ι₁ / ι₂ in under-determined positions)
\newcommand{\marklam}[1][\cdot]{{\markopen #1 \markclose}_{\sim \Rightarrow}}
\newcommand{\markinj}[1][\cdot]{{\markopen #1 \markclose}_{\sim +}}
\newcommand{\markprodsim}[1][\cdot]{{\markopen #1 \markclose}_{\sim \times}}

% Free-variable (unbound) mark ⟨k⟩⇑
\newcommand{\markvar}[1]{{\color{markcolor}\langle} #1 {\color{markcolor}\rangle}\!\Uparrow}

% Red-dashed box around a marked term (error location)
\newcommand{\markbox}[1]{%
  \tikz[baseline=(N.base)]{%
    \node[draw=BrickRed!85, dashed, line width=0.7pt, rounded corners=2pt,
          inner sep=2pt, fill=BrickRed!6, anchor=base] (N) {\strut$#1$};%
  }%
}

% Lighter \sli[OliveGreen] for structural parts of an analysis-slice context
\newcommand{\slist}[1]{%
  \tikz[baseline=(N.base)]{%
    \node[draw=OliveGreen!25, dashed, line width=0.4pt, rounded corners=1.5pt,
          inner sep=1.5pt, fill=OliveGreen!4, anchor=base] (N) {\strut$#1$};%
  }%
}

% Per-mark debugging query block: #1 = υ (syn), #2 = ψ (ana)
\newcommand{\queryblock}[2]{%
  \begin{center}
  \small
  \begin{tabular}{r@{\;}c@{\;}l}
    {\color{NavyBlue}synthesis query}  & $\upsilon\;=$ & $#1$ \\
    {\color{OliveGreen}analysis query} & $\psi\;=$ & $#2$ \\
  \end{tabular}
  \end{center}
}

% Structural variables (from formalism)
\newcommand{\C}{\mathcal{C}}      % expression contexts
\newcommand{\Cs}{c}                % context slices
\newcommand{\p}{d}                 % context typing slices
\newcommand{\F}{\mathcal{F}}      % typing-assumption contexts (endomorphisms)
\newcommand{\f}{f}                 % typing-assumption context slices
\newcommand{\Sb}{\mathcal{S}}     % type-indexed slices

% Analysis slices (cast section): use dutchcal A
\DeclareMathAlphabet{\mathdcal}{U}{dutchcal}{m}{n}
\newcommand{\A}{\mathdcal{A}}

% Type properties (predicate-style)
\newcommand{\ground}[1][\tau]{#1\text{ ground}}
\newcommand{\final}[1][d]{#1\text{ final}}
\newcommand{\val}[1][d]{#1\text{ val}}
\newcommand{\boxedval}[1][d]{#1\text{ boxedval}}
\newcommand{\indet}[1][d]{#1\text{ indet}}

% Contextual substitution [d/u]
\newcommand{\contextualsub}[1][d/u]{\llbracket #1 \rrbracket}

% Polymorphism
\newcommand{\Forall}[2]{\forall #1.\ #2}        % ∀α.τ
\newcommand{\TLam}[2]{\Lambda #1.\ #2}           % Λα.e
\newcommand{\TApp}[2]{#1\ @#2}                   % e @τ
\newcommand{\ForallSlice}[2]{\forall #1.\ #2}    % ∀α.υ
\newcommand{\gapforall}{{\color{teal}\forall\Box.\Box}}

% Labelled tuples
\newcommand{\lbl}[2]{\mathtt{#1} : #2}
\newcommand{\tuplety}[1]{\{#1\}}

% Mechanisation status symbols: ○ yes (green), × no (red), ⊗ partial (orange)
\newcommand{\yo}{{\color{OliveGreen}$\bigcirc$}\;}
\newcommand{\no}{{\color{red}$\times$}\;}
\DeclareRobustCommand{\po}{{\color{orange}$\mathbin{\ooalign{$\bigcirc$\cr\hidewidth$\times$\hidewidth}}$}\;}

% Theorem style: "○ Theorem 4.1 (Name)."
\newtheoremstyle{mechplain}%
  {\topsep}{\topsep}%
  {\itshape}{}%
  {\bfseries}{.}{ }%
  {\thmname{#1} \thmnumber{#2}\thmnote{\;\normalfont(#3)}}%

\theoremstyle{mechplain}
\newtheorem{theorem}{Theorem}[section]
\newtheorem{lemma}[theorem]{Lemma}
\newtheorem{proposition}[theorem]{Proposition}
\newtheorem{corollary}[theorem]{Corollary}
\newtheorem{conjecture}[theorem]{Conjecture}
\newtheorem{counterexample}[theorem]{Counterexample}

\theoremstyle{definition}
\newtheorem{definition}[theorem]{Definition}
\newtheorem{example}[theorem]{Example}

\theoremstyle{remark}
\newtheorem{remark}[theorem]{Remark}

\pagestyle{fancy}
\fancyhf{}
\fancyhead[L]{\small\nouppercase{\leftmark}}
\fancyhead[R]{\thepage}
\renewcommand{\headrulewidth}{0.4pt}
\setlength{\headheight}{14pt}

% Margin annotations on the LEFT
\reversemarginpar
\setlength{\marginparwidth}{2.1cm}
\setlength{\marginparsep}{0.15cm}
% \marginnote survives tcolorbox / theorem envs where \marginpar fails
\usepackage{marginnote}

% Section-level: status symbol + Agda file path, aligned with section heading
\newcommand{\mech}[2]{\marginpar{\raggedleft%
  {\Large#1}\,{\tiny\texttt{#2}}}}
\newcommand{\mechyes}[1]{\mech{{\color{OliveGreen}$\bigcirc$}}{#1}}
\newcommand{\mechno}[1]{\mech{{\color{red}$\times$}}{#1}}
\newcommand{\mechpartial}[1]{\mech{{\color{orange}$\mathbin{\ooalign{$\bigcirc$\cr\hidewidth$\times$\hidewidth}}$}}{#1}}
% \mechfile path; optional first arg = extra vertical offset (default lifts to section heading)
\newcommand{\mechfile}[2][-2\baselineskip]{\marginnote[\raggedleft\tiny\color{black!55}\texttt{#2}]{\raggedleft\tiny\color{black!55}\texttt{#2}}[#1]}

% Box-internal margin annotation (place inside tcolorbox after \begin{...box})
\newcommand{\mechbox}[1]{\marginnote[\raggedleft\tiny\color{black!55}\texttt{#1}]{}}

% Inline status symbols for theorem/corollary headers
\newcommand{\iy}{$\bigcirc$\;}
\newcommand{\ino}{$\times$\;}
\newcommand{\ipart}{$\bigcirc\!\times$\;}

\newtcolorbox{recipe}[1]{
  colback=blue!3,
  colframe=blue!40,
  fonttitle=\bfseries,
  title={Recipe: #1},
  sharp corners,
  boxrule=0.5pt
}

\newtcolorbox{originalerror}[1][]{
  colback=red!3,
  colframe=red!40,
  fonttitle=\bfseries,
  title={Error (Original)#1},
  sharp corners,
  boxrule=0.5pt
}

\newtcolorbox{majortheorembox}{
  colback=green!3,
  colframe=OliveGreen!60,
  sharp corners,
  boxrule=0.5pt,
  breakable,
  enhanced
}

% Master-theorem sub-numbering: produces "Theorem N.M.k"
\newcounter{subthm}
\newcommand{\masterthmnum}{??}
\renewcommand{\thesubthm}{\masterthmnum.\arabic{subthm}}
\newcommand{\mastercases}{%
  \xdef\masterthmnum{\thetheorem}%
  \setcounter{subthm}{0}%
}
\newenvironment{thmcase}[1][]{%
  \refstepcounter{subthm}%
  \begin{majortheorembox}%
  \par\noindent{\bfseries Theorem~\thesubthm}%
  \ifx\\#1\\\else\ (\textit{#1})\fi%
  {\bfseries.}\itshape\quad\ignorespaces
}{%
  \par\end{majortheorembox}%
}

\newtcolorbox{counterexamplebox}{
  colback=orange!3,
  colframe=orange!50,
  sharp corners,
  boxrule=0.5pt,
  breakable,
  enhanced
}

% Inline mini-icons for the \sectionlegend
\newcommand{\thmboxicon}{%
  \tikz[baseline=-0.6ex]{%
    \node[draw=OliveGreen!60, fill=green!3, line width=0.5pt,
          inner sep=0pt, minimum width=10pt, minimum height=6pt] {};%
  }%
}
\newcommand{\ctrboxicon}{%
  \tikz[baseline=-0.6ex]{%
    \node[draw=orange!50, fill=orange!3, line width=0.5pt,
          inner sep=0pt, minimum width=10pt, minimum height=6pt] {};%
  }%
}

% Section-top legend covering mechanisation symbols and callout boxes
\newcommand{\sectionlegend}{%
  \noindent\textit{Legend:}\ \yo fully mechanised, \po mechanised with postulates, \no not mechanised, \thmboxicon\ important theorem, \ctrboxicon\ counterexample.
  \par\medskip
}

% Multi-slice highlighting: A bold red, B/C/D underlines in distinct colours
\newcommand{\sliceA}[1]{\textbf{\underline{\textcolor{BrickRed}{#1}}}}
\newcommand{\sliceB}[1]{\underline{\textcolor{NavyBlue}{#1}}}
\newcommand{\sliceC}[1]{\underline{\textcolor{OliveGreen}{#1}}}
\newcommand{\sliceD}[1]{\underline{\textcolor{Mulberry}{#1}}}

% Highlighted-box slice notation. #1 (opt) = colour (default NavyBlue), #2 = math
\newcommand{\sli}[2][NavyBlue]{%
  \tikz[baseline=(N.base)]{%
    \node[draw=#1!60, dashed, line width=0.4pt, rounded corners=1.5pt,
          inner sep=1.5pt, fill=#1!12, anchor=base] (N) {\strut$#2$};%
  }%
}

% Removable-region marker (red outline over a slice highlight)
\newcommand{\slr}[1]{%
  \tikz[baseline=(N.base)]{%
    \node[draw=BrickRed!80, line width=0.8pt, rounded corners=2pt,
          inner sep=3pt, fill=none, anchor=base] (N) {\strut$#1$};%
  }%
}

% Full slice: assumption box stacked above the sliced term
\newcommand{\fullslice}[2]{%
  \begin{tikzpicture}[baseline=(t.base), every node/.style={inner sep=3pt}]
    \node[draw, line width=0.9pt, rounded corners,
          label={[font=\scriptsize\itshape]above:assumptions}] (a) {$#1$};
    \node[below=6pt of a] (t) {$#2$};
  \end{tikzpicture}%
}

\newcommand{\gapbox}{\boxdot}     % dotted square (gap in expression context)

% Word count: generated by texcount during nix build
% Local: texcount -inc -sum=1,0,0,0,0,0,0 -0 PolymorphicTypeSlicing.tex > wordcount.aux
\newcommand{\wordcount}{\IfFileExists{wordcount.aux}{\input{wordcount.aux}}{N/A}}

%%% Local Variables:
%%% mode: LaTeX
%%% TeX-master: "PolymorphicTypeSlicing"
%%% End:
